\documentclass[11pt]{article}
\usepackage[margin=1in]{geometry}
\usepackage{amsmath,amssymb,booktabs,array,longtable}
\usepackage[hidelinks]{hyperref}

\title{Cumulative Notes on Wage Interaction Models and Firm Effects}
\author{}
\date{\today}

\newcommand{\M}{\mathcal{M}}
\newcommand{\A}{\mathcal{A}}

\begin{document}
\maketitle

\section{Shimer--Smith Bellman Derivation}

\paragraph{Source convention.}
Page citations to Shimer--Smith refer to the printed journal page numbers in
the JSTOR PDF. Page citations to the lecture notes refer to the slide/PDF page
number in the 325-page lecture deck. Formulas marked as ``derived'' are not
written in exactly that form in the source; they are algebraic or first-order
continuous-time implications of the cited source formulas.

\subsection{Notation}

To avoid overloading notation, I write the capitalized object $S(x\mid y)$ for
the personal capital surplus and the lowercase object $s(x,y)$ for the flow match
surplus. This is the same distinction Shimer--Smith use in substance when they
move from equations (4) to (6). I also write the lecture-slide discount rate as
$\rho_L$ when discussing the generic HJB, because Shimer--Smith use $\rho$ for
the meeting-rate parameter. Similarly, I write the lecture-slide instantaneous
payoff as $w_L$ because Shimer--Smith use $w(x)=rW(x)$ for the reservation
wage. In Shimer--Smith, the lecture's $\rho_L$ is $r$.

\begin{center}
\begin{tabular}{>{\raggedright\arraybackslash}p{0.22\linewidth}
                >{\raggedright\arraybackslash}p{0.70\linewidth}}
\toprule
Symbol & Meaning \\
\midrule
$x,y\in[0,1]$ & Agent types. Source: Shimer--Smith, p. 346. \\
$\ell(x)$ & Exogenous population density of type $x$. Source: Shimer--Smith, p. 346. \\
$u(x)$ & Steady-state density of unmatched agents of type $x$. Source: Shimer--Smith, p. 347. \\
$\rho$ & Shimer--Smith meeting-rate parameter. An unmatched type $x$ meets unmatched types in a set $Y$ at rate $\rho\int_Y u(y)\,dy$. Source: Shimer--Smith, p. 347. \\
$\delta$ & Shimer--Smith Poisson separation/dissolution rate of an existing match. This is distinct from the discount rate. Source: Shimer--Smith, pp. 346--348. \\
$r$ & Shimer--Smith discount/interest rate. Source: Shimer--Smith, p. 347. \\
$\rho_L$ & Lecture-slide discount rate in $\rho_L Vdt=w_L\,dt+E[dV]$. The slide writes this as $\rho$; I relabel it to avoid conflict with Shimer--Smith's meeting rate. Source: lecture deck, p. 193. \\
$w_L$ & Lecture-slide instantaneous payoff. The slide writes this as $w(x,u,t)$; I relabel it to avoid conflict with Shimer--Smith's $w(x)=rW(x)$. Source: lecture deck, p. 193. \\
$g_i$ & Instantaneous payoff in state $i$ in my derived jump-process translation of the lecture HJB. In Shimer--Smith, $g_{U_x}=0$ and $g_{M_{xy}}=\pi(x\mid y)$. Derived from Shimer--Smith, pp. 347--348, and lecture deck, p. 193. \\
$\lambda_{ij}$ & Poisson transition rate from state $i$ to state $j$ in my derived jump-process translation. Derived from the Poisson event descriptions in Shimer--Smith, pp. 347--348. \\
$f(x,y)$ & Flow output produced by a matched pair $(x,y)$. Source: Shimer--Smith, pp. 346--347. \\
$\pi(x\mid y)$ & Endogenous flow payoff received by type $x$ while matched with type $y$. Feasibility requires $\pi(x\mid y)+\pi(y\mid x)=f(x,y)$. Source: Shimer--Smith, p. 347. \\
$W(x)$ & Present discounted value of being unmatched for type $x$. Source: Shimer--Smith, p. 348. \\
$W(x\mid y)$ & Present discounted value for type $x$ while matched with type $y$. Source: Shimer--Smith, p. 348. \\
$S(x\mid y)$ & Personal capital surplus from the match: $S(x\mid y)\equiv W(x\mid y)-W(x)$. Source: Shimer--Smith, p. 348. \\
$w(x)$ & Flow value of being unmatched, defined by $w(x)\equiv rW(x)$. Shimer--Smith call this a reservation wage. Source: Shimer--Smith, p. 349. \\
$s(x,y)$ & Flow match surplus, $s(x,y)\equiv f(x,y)-w(x)-w(y)$. Source: Shimer--Smith, p. 349. \\
$A(x)$ & Set of types type $x$ is willing to accept. Source: Shimer--Smith, p. 347. \\
$\M(x)$ & Mutually agreeable matching set: $y\in\M(x)$ iff both $x$ and $y$ accept the match. Source: Shimer--Smith, p. 347. \\
$\theta$ & Reduced meeting/friction coefficient in the value equation, $\theta\equiv \rho/[2(r+\delta)]$. Here $\rho$ is Shimer--Smith's meeting rate, not the lecture discount rate. Source: Shimer--Smith, p. 349. \\
\bottomrule
\end{tabular}
\end{center}

\subsection{From the lecture-slide HJB to a Poisson-jump Bellman equation}

The lecture slides write the continuous-time Bellman equation as
\[
\rho_L V(x,t)\,dt
=
\max_u\left\{w_L(x,u,t)\,dt+E[dV]\right\}.
\]
Source: lecture deck, p. 193, with $\rho$ relabeled as $\rho_L$ and the
instantaneous payoff $w$ relabeled as $w_L$.
For a diffusion state, Ito's Lemma gives
\[
E[dV]
=
\left[
V_t+V_xa+\frac{1}{2}V_{xx}b^2
\right]dt.
\]
Source: lecture deck, p. 194. Shimer--Smith do not have a diffusion state
with drift $a$ and volatility $b$. Their value changes because Poisson events move
the agent between discrete economic states: unmatched, matched with $y$, and
back to unmatched after separation.

For a stationary continuous-time jump process, the same lecture Bellman becomes
\[
\rho_L V_i\,dt
=
g_i\,dt+E[dV_i],
\]
where $i$ denotes the current state and $g_i$ is the current flow payoff. If state
$i$ jumps to states $j$ at Poisson rates $\lambda_{ij}$, then
\[
E[dV_i]
=
\sum_j \lambda_{ij}\bigl(V_j-V_i\bigr)\,dt.
\]
Derived: this replaces the Ito expression for $E[dV]$ on lecture-deck p. 194
with the first-order expectation of Poisson jumps. It is not written in this
generic notation in Shimer--Smith.
So the jump-process version of the lecture equation is
\[
\rho_L V_i
=
g_i+\sum_j \lambda_{ij}\bigl(V_j-V_i\bigr).
\]
Derived from the lecture HJB on p. 193 and the preceding Poisson-jump
expectation.
In Shimer--Smith, set $\rho_L=r$. The entire issue is therefore to identify, for
each state, the current flow payoff $g_i$ and the expected capital gain or loss
$\sum_j \lambda_{ij}(V_j-V_i)$.

\subsection{Unmatched Bellman equation}

For an unmatched type $x$, the flow payoff is normalized to zero. The only
payoff-relevant event over a short interval $dt$ is meeting an acceptable unmatched
type $y\in\M(x)$. This occurs at rate $\rho u(y)\,dy$, and when it occurs the value
jumps from $W(x)$ to $W(x\mid y)$, a capital gain
\[
S(x\mid y)=W(x\mid y)-W(x).
\]
Source: Shimer--Smith, p. 348, for the definition of personal surplus.
Using the lecture-slide Bellman with state $i=U_x$, value $V_i=W(x)$,
instantaneous payoff $g_i=0$, and transition rate
$\lambda_{U_x,M_{xy}}=\rho u(y)\,dy$ for $y\in\M(x)$, we have
\[
E[dW(x)]
=
\rho\int_{\M(x)}\bigl(W(x\mid y)-W(x)\bigr)u(y)\,dy\,dt
=
\rho\int_{\M(x)}S(x\mid y)u(y)\,dy\,dt.
\]
Derived from the meeting technology and matching-set definition on
Shimer--Smith p. 347, plus the surplus definition on p. 348.
Therefore the lecture equation $rVdt=g\,dt+E[dV]$ gives
\[
rW(x)\,dt
=0\cdot dt+\rho\int_{\M(x)}S(x\mid y)u(y)\,dy\,dt.
\]
Dividing by $dt$ gives Shimer--Smith's unmatched Bellman equation:
\begin{equation}
rW(x)=\rho\int_{\M(x)}S(x\mid y)u(y)\,dy.
\tag{SS-2}
\end{equation}
Source: Shimer--Smith equation (2), p. 348. The preceding lines derive it from
the lecture HJB.

\subsection{Matched Bellman equation}

Now suppose type $x$ is currently matched with type $y$. The flow payoff is
$\pi(x\mid y)$, not $w(x)$. The value $w(x)=rW(x)$ is a continuation-value object
attached to the outside option of being unmatched.

Apply the lecture-slide Bellman to the matched state $i=M_{xy}$, whose value is
$V_i=W(x\mid y)$ and whose instantaneous payoff is
$g_i=\pi(x\mid y)$. The only jump is separation, which moves $x$ from
$M_{xy}$ to the unmatched state $U_x$ at Poisson rate $\delta$. Therefore
\[
E[dW(x\mid y)]
=\delta dt\,[W(x)-W(x\mid y)]+o(dt)
=-\delta S(x\mid y)\,dt+o(dt).
\]
Derived from the separation description on Shimer--Smith pp. 346--348 and the
definition $S(x\mid y)=W(x\mid y)-W(x)$ on p. 348.
Substituting this into the lecture equation yields
\[
rW(x\mid y)\,dt
=\pi(x\mid y)\,dt-\delta S(x\mid y)\,dt+o(dt).
\]
Dividing by $dt$ and letting $dt\downarrow0$ gives
\begin{equation}
rW(x\mid y)=\pi(x\mid y)-\delta S(x\mid y).
\tag{SS-3}
\end{equation}
Source: Shimer--Smith equation (3), p. 348. The expected-capital-loss line
above derives the otherwise unstated $-\delta S(x\mid y)$ term from the lecture
HJB.
This is the term the paper states but does not derive: the ``loss from
continuing'' is not an extra flow utility term. It is the expected capital loss from
the Poisson jump out of the match into the unmatched state.

\subsection{Equivalent small-\texorpdfstring{$dt$}{dt} derivation}

The same matched equation can be derived by writing the recursive value directly:
\[
W(x\mid y)
=\pi(x\mid y)dt
 + e^{-rdt}\left[(1-\delta dt)W(x\mid y)+\delta dt\,W(x)\right]
 + o(dt).
\]
Derived small-$dt$ recursion from Shimer--Smith's flow payoff and separation
process on pp. 347--348. This recursive equation is not printed in the paper.
Using $e^{-rdt}=1-rdt+o(dt)$, subtracting $W(x\mid y)$ from both sides, and
dropping terms of order $o(dt)$ gives
\[
0=\pi(x\mid y)dt-rW(x\mid y)dt+\delta dt\,[W(x)-W(x\mid y)].
\]
Rearranging:
\[
rW(x\mid y)=\pi(x\mid y)-\delta\bigl(W(x\mid y)-W(x)\bigr)
=\pi(x\mid y)-\delta S(x\mid y).
\]

\subsection{From Bellman equations to Shimer--Smith's surplus formula}

Under Nash surplus division, Shimer--Smith impose
\[
S(x\mid y)=S(y\mid x).
\]
Source: Shimer--Smith, p. 348.
Using the matched Bellman equation for both agents gives
\[
rW(x\mid y)=\pi(x\mid y)-\delta S(x\mid y),
\qquad
rW(y\mid x)=\pi(y\mid x)-\delta S(y\mid x).
\]
Since $W(x\mid y)=W(x)+S(x\mid y)$, the first equation can be rewritten as
\[
rW(x)+(r+\delta)S(x\mid y)=\pi(x\mid y).
\]
Similarly,
\[
rW(y)+(r+\delta)S(y\mid x)=\pi(y\mid x).
\]
Add these two equations, use feasibility
$\pi(x\mid y)+\pi(y\mid x)=f(x,y)$, and use equal surplus division:
\[
rW(x)+rW(y)+2(r+\delta)S(x\mid y)=f(x,y).
\]
Derived algebraically from Shimer--Smith equations (3) and the feasibility
condition on p. 347.
Therefore
\begin{equation}
S(x\mid y)
=\frac{f(x,y)-rW(x)-rW(y)}{2(r+\delta)}.
\tag{SS-4}
\end{equation}
Source: Shimer--Smith equation (4), p. 348. The preceding lines show the
algebra.
Defining $w(x)=rW(x)$ and $s(x,y)=f(x,y)-w(x)-w(y)$, this is
\[
S(x\mid y)=\frac{s(x,y)}{2(r+\delta)}.
\]
Derived from the definitions on Shimer--Smith p. 349 and equation (4) on
p. 348.
Substituting into the unmatched Bellman equation gives
\[
rW(x)
=\rho\int_{\M(x)}
\frac{f(x,y)-w(x)-w(y)}{2(r+\delta)}u(y)\,dy.
\]
Since $w(x)=rW(x)$,
\begin{equation}
w(x)
=\theta\int_{\M(x)}
\bigl(f(x,y)-w(x)-w(y)\bigr)u(y)\,dy
=\theta\int_{\M(x)}s(x,y)u(y)\,dy,
\qquad
\theta=\frac{\rho}{2(r+\delta)}.
\tag{SS-6}
\end{equation}
Source: Shimer--Smith equation (6), p. 349, derived by substituting equation
(4) into equation (2).

\subsection{Main interpretation}

The paper's equations are internally consistent once we keep three objects
separate. First, $\pi(x\mid y)$ is the current flow payoff while matched. Second,
$W(x)$ and $W(x\mid y)$ are capitalized continuation values. Third, $w(x)=rW(x)$
is a flow-equivalent value of the unmatched outside option, not an additional
matched flow payoff. The term $-\delta S(x\mid y)$ in the matched Bellman
equation is the expected capital loss from separation: with probability
$\delta dt$, the agent loses $W(x\mid y)-W(x)$ over the interval $dt$.

\section{Crippa Appendix H: Cycles, Connectedness, and Identification}

\paragraph{Source convention.}
Page citations in this section refer to the printed page numbers in Crippa
(2025), \emph{Identification, Estimation, and Inference in Two-Sided Interaction
Models}. The relevant source passages are Theorems 1--3 and Assumptions 1--3
on pp. 14--19, and proofs H.1--H.3 on pp. 61--66. Formulas marked as
``derived'' are algebraic unpackings of those cited formulas.

\subsection{Notation}

In this section, the user's ``$\beta$'' is Crippa's Tukey interaction parameter
$\beta_0$.

\begin{center}
\begin{tabular}{>{\raggedright\arraybackslash}p{0.24\linewidth}
                >{\raggedright\arraybackslash}p{0.68\linewidth}}
\toprule
Symbol & Meaning \\
\midrule
$I,J$ & Worker and firm node sets in the bipartite matching network. Source: Crippa, Definition 2, p. 14. \\
$G_{IJ}$ & Bipartite matching network with observed edges. Source: Crippa, Definition 2, p. 14. \\
$(i,j)$ & Observed worker-firm edge/match. Source: Crippa, p. 14. \\
$\theta_{ij}$ & Systematic outcome on edge $(i,j)$. Source: Crippa, proofs H.1--H.2, pp. 61--64. \\
$\alpha_i$ & Worker productivity/effect in the Tukey model. Source: Crippa, H.1--H.2, pp. 61--64. \\
$\psi_j$ & Firm productivity/effect in the Tukey model. Source: Crippa, H.1--H.2, pp. 61--64. \\
$\beta_0$ & Tukey interaction parameter. Source: Crippa, Theorem 1 and H.1, pp. 14, 61--63. \\
$K$ & Half-length of a cycle. A cycle of length $2K$ has $K\geq2$. Source: Crippa, H.1, p. 62. \\
$u_i,v_j$ & H.1 transformed variables $u_i=1+\beta_0\alpha_i$ and $v_j=1+\beta_0\psi_j$. Source: Crippa, H.1, p. 62. \\
$S_r^{\mathrm{odd}},S_r^{\mathrm{even}}$ & Elementary symmetric sums of odd- and even-indexed cycle outcomes. Source: Crippa, H.1, pp. 62--63. \\
$\Delta_r$ & Difference $S_r^{\mathrm{odd}}-S_r^{\mathrm{even}}$. Source: Crippa, H.1, p. 63. \\
$\alpha_i',\psi_j',\theta_{ij}'$ & H.2 transformed variables: $\alpha_i'=1+\beta_0\alpha_i$, $\psi_j'=1+\beta_0\psi_j$, and $\theta_{ij}'=1+\beta_0\theta_{ij}$. Source: Crippa, H.2, p. 64. \\
$a_j,b_j$ & Firm-specific intercept and slope in the BLM model. Source: Crippa, H.3, p. 65. \\
$S_n^I,S_n^J$ & Worker and firm sets reached by the Seed-and-Snowballs construction. Source: Crippa, Definition 3 and H.3, pp. 17--18, 65. \\
\bottomrule
\end{tabular}
\end{center}

\subsection{What H.1 Says About the Cycle Needed for \texorpdfstring{$\beta_0$}{beta0}}

Crippa's Tukey model is
\[
\theta_{ij}=\alpha_i+\psi_j+\beta_0\alpha_i\psi_j .
\]
Source: Crippa, H.1 and H.2, pp. 61, 64.

The main-text Assumption 1 is explicitly a \emph{length-4} informative cycle:
the network contains a cycle $i_1,j_1,i_2,j_2$ with heterogeneous worker and firm
productivities. Source: Crippa, Assumption 1, p. 14. Theorem 1 then says
Assumption 1 is sufficient for identifying $\beta_0$, and, if the network contains
no more than one cycle, it is also necessary. Source: Crippa, Theorem 1, p. 14.

Appendix H.1 explains why the length matters. For a general cycle of length
$2K$ with $K\geq2$, define
\[
u_i:=1+\beta_0\alpha_i,
\qquad
v_j:=1+\beta_0\psi_j.
\]
Then each observed edge satisfies
\[
1+\beta_0\theta_{ij}
=1+\beta_0(\alpha_i+\psi_j+\beta_0\alpha_i\psi_j)
=(1+\beta_0\alpha_i)(1+\beta_0\psi_j)
=u_iv_j.
\]
Source: Crippa, H.1, p. 62.

If the cycle is ordered as
\[
i_1-j_1-i_2-j_2-\cdots-i_K-j_K-i_1,
\]
then traversing the closed cycle gives
\[
\prod_{k=1}^{K}\left(1+\beta_0\theta^{(2k-1)}\right)
=
\prod_{k=1}^{K}u_{i_k}v_{j_k}
=
\prod_{k=1}^{K}u_{i_{k+1}}v_{j_k}
=
\prod_{k=1}^{K}\left(1+\beta_0\theta^{(2k)}\right),
\qquad i_{K+1}=i_1.
\]
Source: Crippa, H.1, p. 62.

This is the exact point where cycle traversal is used for $\beta_0$. The product
collapse requires a closed path: the odd-edge product and even-edge product use
the same multiset of worker factors and the same multiset of firm factors. Without
closing the path, those factors would not cancel/rearrange into the same product.
Derived interpretation of the source equality on Crippa, H.1, p. 62.

Expanding the two products in elementary symmetric sums and subtracting gives
\[
\sum_{r=1}^{K}\left(S_r^{\mathrm{odd}}-S_r^{\mathrm{even}}\right)\beta_0^r=0.
\]
With $\Delta_r=S_r^{\mathrm{odd}}-S_r^{\mathrm{even}}$, this is
\[
\Delta_1\beta_0+\Delta_2\beta_0^2+\cdots+\Delta_K\beta_0^K=0.
\]
Source: Crippa, H.1, p. 63.

Because $\beta_0=0$ is always an algebraic root, H.1 factors it out and focuses
on
\[
\Delta_1+\Delta_2\beta_0+\Delta_3\beta_0^2+\cdots+\Delta_K\beta_0^{K-1}=0.
\]
Source: Crippa, H.1, p. 63.

Thus a single cycle delivers one polynomial restriction:
\[
P(\beta_0)=\beta_0 Q(\beta_0)=0,\qquad \deg Q\le K-1.
\]
The root $\beta=0$ is algebraic and uninformative. After removing it, a 4-cycle
has $K=2$, so $Q$ is linear and the nonzero root is unique under Crippa's
informativeness condition. If $K>2$, $Q$ is higher degree, so one cycle gives a
finite candidate set for $\beta_0$ rather than point identification. Source:
Crippa, H.1, p. 63.

The main text adds the important qualification: if there are multiple longer
cycles, the intersection of their identification sets can reduce to a singleton.
Source: Crippa, p. 15. So the answer to the practical question is:
\[
\text{Crippa's sufficient primitive condition is an informative 4-cycle.}
\]
\[
\text{A longer }2K\text{-cycle gives restrictions, but one such cycle alone does not
generically give a unique }\beta_0.
\]
\[
\text{Multiple longer cycles may still point-identify by intersecting root sets.}
\]

For the special length-4 case, H.1 gives
\[
\beta_0
=-
\frac{\theta_{i_1j_1}+\theta_{i_2j_2}-\theta_{i_1j_2}-\theta_{i_2j_1}}
{\theta_{i_1j_1}\theta_{i_2j_2}-\theta_{i_1j_2}\theta_{i_2j_1}},
\]
with Assumption 1 ensuring the denominator is nonzero. Source: Crippa, H.1,
pp. 63--64.

\subsection{What H.2 Uses Connectedness For}

H.2 assumes $\beta_0$ is already known and proves identification of the Tukey
productivities under connectedness. Source: Crippa, Theorem 2, p. 16; H.2,
p. 64. If $\beta_0\neq0$, H.2 rewrites the Tukey model as
\[
1+\beta_0\theta_{ij}
=(1+\beta_0\alpha_i)(1+\beta_0\psi_j),
\]
and sets
\[
\alpha_i'=1+\beta_0\alpha_i,
\qquad
\psi_j'=1+\beta_0\psi_j,
\qquad
\theta_{ij}'=1+\beta_0\theta_{ij}.
\]
Thus
\[
\theta_{ij}'=\alpha_i'\psi_j'.
\]
Source: Crippa, H.2, p. 64.

Connectedness is enough here because identification can be propagated along any
path. With the normalization $\alpha_{i_0}=0$, we have $\alpha_{i_0}'=1$.
If workers $i$ and $i'$ both match with firm $j$, then
\[
\frac{\theta_{ij}'}{\theta_{i'j}'}
=
\frac{\alpha_i'\psi_j'}{\alpha_{i'}'\psi_j'}
=
\frac{\alpha_i'}{\alpha_{i'}'}.
\]
Source: Crippa, H.2, p. 64.

For a path
\[
i_0-j_1-i_1-j_2-\cdots-j_m-i,
\]
H.2 multiplies these local ratios:
\[
\frac{\alpha_i'}{\alpha_{i_0}'}
=
\frac{\alpha_i'}{\alpha_{i_{m-1}}'}
\cdot
\frac{\alpha_{i_{m-1}}'}{\alpha_{i_{m-2}}'}
\cdots
\frac{\alpha_{i_1}'}{\alpha_{i_0}'}
=
\prod_{k=1}^{m}\frac{\theta_{i_kj_k}'}{\theta_{i_{k-1}j_k}'}.
\]
Source: Crippa, H.2, p. 64.

This is a path propagation argument, not a cycle identification argument for
$\beta_0$. Once $\beta_0$ is known, connectedness lets every worker be reached
from $i_0$, and then each observed edge gives
\[
\psi_j'=\frac{\theta_{ij}'}{\alpha_i'}.
\]
Source: Crippa, H.2, p. 64. Necessity is the standard component-wise rescaling
problem: if $G_{IJ}$ is disconnected, each component can be rescaled separately
without changing products $\alpha_i'\psi_j'$. Source: Crippa, H.2, p. 64.

\subsection{What H.3 Adds for BLM}

H.3 concerns the BLM model, where firm-specific parameters enter as
\[
\theta_{ij}=a_j+b_j\alpha_i.
\]
This displayed form is source from H.3, p. 65. H.3 first notes that two
normalizations are needed. In particular, the transformation
\[
\alpha_i'=c_1\alpha_i+c_2,
\qquad
b_j'=\frac{b_j}{c_1},
\qquad
a_j'=a_j-\frac{b_jc_2}{c_1}
\]
is observationally equivalent for $c_1\neq0$. Source: Crippa, H.3, p. 65.

The sufficiency proof uses the Seed-and-Snowballs construction. If firm
parameters $(a_j,b_j)$ are known for firms already in the snowball, then a
worker connected to one such firm is identified by
\[
\alpha_i=\frac{\theta_{ij}-a_j}{b_j}.
\]
Source: Crippa, H.3, p. 65.

Conversely, if a new firm is connected to two already-known workers $i,i'$ with
distinct productivities, then
\[
\begin{cases}
\theta_{ij}=a_j+b_j\alpha_i,\\
\theta_{i'j}=a_j+b_j\alpha_{i'},
\end{cases}
\qquad\Longleftrightarrow\qquad
\begin{cases}
b_j=\dfrac{\theta_{i'j}-\theta_{ij}}{\alpha_{i'}-\alpha_i},\\[1.0em]
a_j=\theta_{ij}-\dfrac{\theta_{i'j}-\theta_{ij}}{\alpha_{i'}-\alpha_i}\alpha_i.
\end{cases}
\]
Source: Crippa, H.3, p. 65.

This is why BLM needs more than connectedness. A firm-specific slope $b_j$ cannot
be identified from a single already-known worker; two distinct worker
productivities are needed. Seed-and-Snowballs imposes exactly this iterative
structure: firms bring in workers, and firms newly added must connect to at
least two already-reached workers. Source: Crippa, Definition 3, pp. 17--18;
Assumption 3 and discussion, p. 19; H.3, p. 65.

H.3's necessity argument shows what fails otherwise. If the graph is connected
but the snowball stalls, an outer firm can have at most one bridge worker in the
reached worker set. Crippa then constructs a one-parameter family indexed by
$\lambda$:
\[
b_{j_\lambda}^{\lambda}=\lambda b_{j_\lambda},
\qquad
a_{j_\lambda}^{\lambda}
=\theta_{i_\lambda j_\lambda}
-b_{j_\lambda}^{\lambda}\alpha_{i_\lambda},
\]
which preserves the bridged observed outcome while changing the firm
parameters. Source: Crippa, H.3, p. 66. Propagating this redefinition through the
remaining outer nodes yields infinitely many distinct parameter vectors
consistent with the same data. Source: Crippa, H.3, p. 66.

\subsection{Clean Takeaways for Our Use}

\begin{enumerate}
\item For Tukey's $\beta_0$, Crippa's theorem uses an informative 4-cycle as the
clean sufficient condition. A longer $2K$-cycle gives a degree-$(K-1)$ equation
after removing the trivial zero root, so one longer cycle alone gives a finite set
of candidate values, not guaranteed point identification. Multiple longer cycles
can identify if their candidate sets intersect in one value. Source: Crippa,
pp. 15, 62--63.
\item The cycle traversal is used specifically to identify/restrict $\beta_0$:
the product over odd cycle edges equals the product over even cycle edges because
the closed traversal contains the same worker and firm multiplicative factors on
both sides. Source: Crippa, H.1, p. 62.
\item Once $\beta_0$ is known, connectedness is enough for $\alpha$ and $\psi$ in
the Tukey model. H.2 propagates ratios along paths; no cycle is needed at this
stage. Source: Crippa, H.2, p. 64.
\item BLM requires the stronger Seed-and-Snowballs condition because each new
firm has both an intercept and a slope. Identifying those two firm parameters
requires two already-identified workers with distinct productivities. Source:
Crippa, H.3, p. 65.
\end{enumerate}

\section{Card--Heining--Kline (2013): One-Page Empirical Summary}

\paragraph{Source convention.}
Page citations in this section refer to the NBER working-paper PDF
\emph{Workplace Heterogeneity and the Rise of West German Wage Inequality}
(NBER Working Paper 18522). Formulas and numerical claims are either directly
from the paper or algebraic restatements of cited source equations.

\subsection{Notation}

\begin{center}
\begin{tabular}{>{\raggedright\arraybackslash}p{0.24\linewidth}
                >{\raggedright\arraybackslash}p{0.68\linewidth}}
\toprule
Symbol & Meaning \\
\midrule
$y_{it}$ & Log daily wage of worker $i$ in year $t$. Source: CHK, p. 16. \\
$\alpha_i$ & Time-invariant worker component of pay. Source: CHK, p. 16. \\
$\psi_{J(i,t)}$ & Establishment component for worker $i$'s establishment in year $t$. Source: CHK, p. 16. \\
$x_{it}'\beta$ & Time-varying observable component, including year effects and age/education interactions. Source: CHK, pp. 16--17. \\
$\eta_{iJ(i,t)}$ & Time-invariant match component for worker $i$ at establishment $J(i,t)$. Source: CHK, pp. 16--17. \\
$\zeta_{it}$ & Unit-root component of worker wages. Source: CHK, p. 17. \\
$\varepsilon_{it}$ & Transitory error. Source: CHK, p. 17. \\
$r_{it}$ & Composite residual in the AKM estimating equation. Source: CHK, p. 17. \\
$D,F,X$ & Worker, establishment, and covariate design matrices. Source: CHK, pp. 17--18. \\
\bottomrule
\end{tabular}
\end{center}

\subsection{Summary}

CHK study whether rising West German wage inequality from 1985 to 2009 can be
explained by workplace heterogeneity. Their core empirical model is
\[
y_{it}
=
\alpha_i+\psi_{J(i,t)}+x_{it}'\beta+\eta_{iJ(i,t)}+\zeta_{it}+\varepsilon_{it}
=
\alpha_i+\psi_{J(i,t)}+x_{it}'\beta+r_{it}.
\]
Source: CHK, equation (1), pp. 16--17. In matrix form,
\[
y=D\alpha+F\psi+X\beta+r.
\]
Source: CHK, equation (2), p. 17. The identifying orthogonality conditions are
\[
E[D'r]=0,\qquad E[F'r]=0,\qquad E[X'r]=0.
\]
Source: CHK, equation (4), p. 18. The main empirical concern is whether the
composite error is orthogonal to establishment assignment, which CHK discuss as
an exogenous mobility condition. Source: CHK, pp. 18--20.

The paper uses German administrative data from the IAB/IEB for full-time West
German men aged 20--60, split into four intervals over 1985--2009. Source: CHK,
pp. 4, 8--11. The data are huge: in each interval CHK have about 85--90 million
person-year observations and about 17 million workers; the largest connected set
contains about 97 percent of person-year observations and about 95 percent of
workers. Source: CHK, p. 21.

The headline empirical result is that AKM-style worker and establishment effects
explain most of the rise in West German wage inequality. CHK report that the
standard deviation of log wages rises from 0.38 in 1985--1991 to 0.51 in
2002--2009. Source: CHK, p. 21. Their estimated worker and establishment effects
both become more dispersed, and the correlation between them rises from 0.03 to
0.25. Source: CHK, p. 22. The residual component is small and rises only modestly:
of the 13 log-point rise in the standard deviation of wages, only 1.5 log points
remain unexplained. Source: CHK, p. 22.

Their variance decomposition is the most useful one-line takeaway for us. Between
the first and fourth intervals, the variance of worker effects rises from 0.084 to
0.127, accounting for about 40 percent of the rise in wage variance; the variance of
establishment effects rises from 0.025 to 0.053, contributing about 25 percent; and
the covariance term rises from 0.003 to 0.041, adding about 34 percent. Source:
CHK, p. 25. CHK interpret this as evidence that rising inequality reflects three
forces: more dispersed worker components, more dispersed establishment premia,
and stronger assortative matching between high-wage workers and high-wage
establishments. Source: CHK, abstract p. 1 and p. 25.

CHK also run diagnostics for match effects and endogenous mobility. A saturated
match-effect model fits statistically better, but the gain is small: the implied
standard deviation of match effects rises only from about 0.060 to 0.075. Source:
CHK, pp. 20, 23. They also examine residuals by worker-effect and establishment-
effect deciles; mean residuals in the 100 crossed cells are uniformly below 1 percent
in magnitude, though the low-worker/low-establishment cells show small systematic
departures from additivity. Source: CHK, p. 23. Event studies show roughly
symmetric wage gains and losses for movers between high- and low-wage
establishments, and little evidence of pre-move wage dips or blips. Source: CHK,
pp. 14--16, 23--24. Their conclusion is therefore not that the AKM model is exact,
but that it is a good first-order empirical approximation for their setting.

\section{Sorkin--Warwar: Why the Event-Study Test Can Be Low Power}

\paragraph{Source convention.}
Page citations in this section refer to the PDF \emph{Paired Movers} by Isaac
Sorkin and Lucas Warwar, March 2026. The event-study discussion is mainly on
pp. 31--34, with the broader paper summary on pp. 1--4.

\subsection{Notation}

\begin{center}
\begin{tabular}{>{\raggedright\arraybackslash}p{0.24\linewidth}
                >{\raggedright\arraybackslash}p{0.68\linewidth}}
\toprule
Symbol & Meaning \\
\midrule
$A,B$ & Origin and destination firms for paired movers. Source: SW, p. 1. \\
$h_i$ & Worker effect in SW's search model. Source: SW, p. 27. \\
$p_j$ & Firm effect in SW's search model. Source: SW, p. 27. \\
$m_{ij}$ & Match effect between worker $i$ and firm $j$. Source: SW, p. 27. \\
EE & Employer-to-employer move. Source: SW, pp. 30--32. \\
EUE & Employer-to-unemployment-to-employer move. Source: SW, pp. 30--32. \\
$\lambda_1$ & On-the-job offer arrival parameter; varying it changes the share of EE moves in the simulations. Source: SW, pp. 28, 30--32. \\
Event-study coefficient & Regression/IV coefficient of earnings changes on changes in firm effects. Source: SW, pp. 31--32. \\
\bottomrule
\end{tabular}
\end{center}

\subsection{What the Diagnostic Is Testing}

Sorkin--Warwar describe the CHK event-study test as asking whether changes in
firm effects predict within-worker earnings changes. Source: SW, p. 31. In the
AKM-style interpretation, the diagnostic is reassuring if a worker moving from a
low firm effect to a high firm effect gains approximately the firm-effect difference,
and a worker moving in the reverse direction loses approximately the same amount.
This is why CHK emphasize symmetry of gains/losses and absence of pre-move dips.
Source: CHK, pp. 14--16, 18--20, 23--24.

The link to the no-match-effect/parallel-trends intuition is: if mobility is not
driven by unobserved match-specific wage components, then the change in the
estimated firm component should forecast the change in the worker's wage. In
Kline's edge-effect language, AKM imposes restrictions on edge effects: the effect of
walking around a cycle should sum to zero, so firm effects summarize edge wage
changes without residual cycle effects. Source: Kline, pp. 18--19. SW's point is not
that the event-study diagnostic is meaningless; it is that passing it does not rule out
large match effects or endogenous mobility. Source: SW, pp. 31--34.

\subsection{SW's Low-Power Mechanism}

SW estimate a search model in which a worker's payoff at firm $j$ is
\[
h_i+p_j+m_{ij}.
\]
Source: SW, p. 27. The worker sees both firm and match components when deciding
whether to accept offers, so the model violates both additive separability and
exogenous mobility. Source: SW, p. 27.

The surprising result is that this endogenous-mobility model still passes the
pooled event-study test. In simulated data, SW report an event-study coefficient of
0.99 in the full sample, even though 41 percent of moves are EE moves. Source: SW,
p. 32. Split by move type, the diagnostic behaves as expected: the coefficient is
0.57 for EE moves and 0.99 for EUE moves. Source: SW, pp. 31--32. Thus, the test
does detect the problem in the EE subsample, but the pooled sample hides it.

The reason is a pooling/composition effect. SW explain that OLS pooled coefficients
combine subsamples using relative sample sizes, within-group variances, and
between-group differences in sample means. Source: SW, p. 32. In their model, EE
moves have larger increases in firm effects and earnings than EUE moves, because
endogenous movers tend to accept offers with favorable match and firm components.
This creates a between-sample slope above one, which pushes the pooled coefficient
back toward one even when within-subsample coefficients are below one. Source: SW,
pp. 32--34.

SW show this mechanism is strong. As they vary $\lambda_1$ to increase the share
of EE moves, both the EE and EUE event-study coefficients fall, but the pooled
coefficient remains above 0.90 and can exceed the coefficient in either subsample.
Source: SW, p. 32. In the data, the pooled Brazilian event-study coefficient is
about 1.09, but when moves are split by whether firm effects increase or decrease,
the coefficients are 1.28 and 1.50; SW interpret this as evidence that pooling makes
the event-study test look better than the relevant subsamples. Source: SW, pp.
33--34.

\subsection{Takeaway}

The standard event-study diagnostic is best read as a test of whether realized mover
wage changes line up with estimated firm-effect changes in the pooled mover sample.
It is not a high-power test against match-driven endogenous mobility. SW's
mechanism is that endogenous and more exogenous moves can have different sample
means; when pooled, the between-sample slope can offset within-subsample failures.
So the diagnostic can pass even when a large part of the data violates the
no-match-effect/exogenous-mobility logic. Source: SW, pp. 31--35.

\section{Crippa's BLM Model, Real BLM, and LLMR}

\paragraph{Source convention.}
Page citations to Crippa refer to printed page numbers in Crippa (2025). Page
citations to Bonhomme--Lamadon--Manresa refer to printed Econometrica pages in
BLM (2019). Page citations to Lamadon--Lise--Meghir--Robin refer to the PDF
pages of the December 30, 2025 draft.

\subsection{Notation}

\begin{center}
\begin{tabular}{>{\raggedright\arraybackslash}p{0.24\linewidth}
                >{\raggedright\arraybackslash}p{0.68\linewidth}}
\toprule
Symbol & Meaning \\
\midrule
$\theta_{ij}$ & Systematic outcome for worker $i$ matched to firm $j$ in Crippa's bipartite interaction framework. Source: Crippa, pp. 10, 65. \\
$\alpha_i$ & Worker scalar productivity/type in Crippa's BLM model. Source: Crippa, pp. 10, 65. \\
$a_j,b_j$ & Firm-specific intercept and slope in Crippa's BLM model. Source: Crippa, pp. 10, 65. \\
$k(j)$ or $k_{it}$ & Firm class in BLM. Source: BLM, pp. 702--703. \\
$K$ & Number of BLM firm classes. Source: BLM, pp. 702--703, 712. \\
$\xi_i$ & Worker type in BLM. Source: BLM, p. 703. \\
$Y_{it}$ & Log earnings in BLM. Source: BLM, p. 703. \\
$X_{it}$ & Observed worker covariates in BLM. Source: BLM, p. 703. \\
$m_{it}$ & Mobility indicator in BLM. Source: BLM, p. 703. \\
$a_t(k),b_t(k)$ & Firm-class intercept and slope in BLM's simple interactive regression example. Source: BLM, p. 704. \\
$x,y$ & Worker and firm types in LLMR. Source: LLMR, pp. 6--8. \\
$f_{xy},c_{xy}$ & Match production and disutility/amenity in LLMR. Source: LLMR, pp. 7--9. \\
$R,S_{xy}$ & Worker surplus and maximum feasible worker surplus in LLMR. Source: LLMR, p. 9. \\
\bottomrule
\end{tabular}
\end{center}

\subsection{Crippa's Stripped-Down BLM}

Crippa defines the BLM model as
\[
\theta_{ij}=a_j+b_j\alpha_i.
\]
Source: Crippa, Section 2.3.1 and H.3, pp. 10, 65. The economic content is clean:
each firm has an intercept $a_j$ and a slope $b_j$ multiplying the worker scalar
productivity $\alpha_i$. The slope is the firm-specific complementarity parameter:
firms can differ not only in level productivity, but also in how strongly worker
productivity matters at that firm. Source: Crippa, p. 10.

Crippa explicitly says this model is motivated by Bonhomme et al. (2019), but that
there is a difference in granularity. In BLM, the functional form is embedded in a
grouped setting where all firms and workers within a group share latent
characteristics. In Crippa's bipartite-interaction framework, the BLM model is the
case where each group consists of a single worker or a single firm. Source: Crippa,
p. 10.

This simplification is useful because it makes identification algebra transparent. In
H.3, if $(a_j,b_j)$ are known for a firm and $b_j\neq0$, then a linked worker's
productivity is recovered by
\[
\alpha_i=\frac{\theta_{ij}-a_j}{b_j}.
\]
If a new firm is linked to two already-identified workers $i,i'$ with distinct
productivities, then
\[
\begin{cases}
\theta_{ij}=a_j+b_j\alpha_i,\\
\theta_{i'j}=a_j+b_j\alpha_{i'},
\end{cases}
\qquad\Longleftrightarrow\qquad
\begin{cases}
b_j=\dfrac{\theta_{i'j}-\theta_{ij}}{\alpha_{i'}-\alpha_i},\\[1.0em]
a_j=\theta_{ij}-\dfrac{\theta_{i'j}-\theta_{ij}}{\alpha_{i'}-\alpha_i}\alpha_i.
\end{cases}
\]
Source: Crippa, H.3, p. 65. This is why the Seed-and-Snowballs graph condition is
needed: a new firm needs two already-known workers, not just one. Source: Crippa,
pp. 18--19, 65.

\subsection{What Real BLM Adds}

BLM is not just Crippa's deterministic node-level equation with different notation.
It adds four layers: grouping, distributions, estimation, and dynamics. Source: BLM
abstract, p. 699.

\begin{enumerate}
\item \textbf{Grouping.} BLM use firm classes $k=k(j)\in\{1,\ldots,K\}$ rather than
one parameter per firm. This reduces the incidental-parameter problem created by
few movers per firm. Source: BLM, pp. 700--703. Classes are estimated first by
k-means on earnings distributions, and the model is then estimated conditional on
those classes. Source: BLM, pp. 701, 711--712.

\item \textbf{Distributional worker heterogeneity.} Workers have latent type $\xi_i$,
and BLM recover type--class earnings distributions and type proportions within firm
classes, not just conditional means. Source: BLM, p. 703. In the static model,
mobility may depend on worker type and firm classes, but not directly on earnings;
post-move earnings depend on worker type and the new class, not on the previous
class or previous earnings. Source: BLM, p. 703.
\end{enumerate}

BLM's simple interactive regression example is
\[
Y_{it}=a_t(k_{it})+b_t(k_{it})\xi_i+X_{it}'c_t+\varepsilon_{it}.
\]
Source: BLM, p. 704. This is the direct ancestor of Crippa's
$\theta_{ij}=a_j+b_j\alpha_i$, but BLM's version has calendar-time dependence,
covariates, firm classes rather than individual firms, latent worker types, and an
error/distributional layer.

\begin{enumerate}
\setcounter{enumi}{2}
\item \textbf{Mixture identification and estimation.} With two periods, BLM identify
type- and class-specific earnings distributions and worker-type proportions, up to
latent-type labels, using movers, firm-class cycles, and rank/linear-independence
conditions. Source: BLM, pp. 706--709.

\item \textbf{Dynamics.} BLM also study a dynamic model where mobility may depend
on current earnings and post-move earnings may depend on previous earnings and
previous firm class through a first-order Markov structure. Source: BLM, pp. 704--706.
Their dynamic identification result uses four periods. Source: BLM, pp. 700, 710--711.
\end{enumerate}

\subsection{What Crippa Drops, and Whether It Is a Problem}

\begin{enumerate}
\item \textbf{Grouping/dimension reduction.} Crippa drops BLM's firm-class
dimension reduction and works at the individual node level. This is excellent for
transparent graph identification, but it is empirically much more demanding. Crippa
himself says that BLM's additional flexibility is costly: the full node-level BLM
condition requires a rich network, and in labor data many firms are not even in cycles;
he concludes that point identification is challenging unless one imposes additional
dimension-reduction restrictions such as BLM. Source: Crippa, p. 19.

\item \textbf{Distributional/mixture structure.} Crippa studies a systematic outcome
$\theta_{ij}$; BLM estimate earnings distributions, worker-type mixtures, and type
composition by firm class. This is not a problem if our paper only needs a clean wage
or production equation. It would be a problem if we tried to claim we inherit BLM's
estimation strategy or distributional interpretation.

\item \textbf{Covariates, time, noise, and likelihood.} Crippa abstracts from
$X_{it}$, $t$-specific wage functions, residual distributions, and maximum-likelihood
estimation. BLM include these because their goal is empirical estimation in panel
data. Source: BLM, pp. 703--704, 713--714. For our modeling/nesting exercise, this
is probably harmless; for empirical implementation, it is not.

\item \textbf{Mobility assumptions.} Crippa's BLM identification result is a graph
condition for recovering deterministic primitives from observed links. BLM's static
model imposes restrictions on mobility and serial dependence, while the dynamic model
relaxes strict exogeneity by allowing earnings to affect mobility. Source: BLM,
pp. 703--706. So Crippa is not carrying over BLM's mobility model; he is isolating
the wage-interaction form.

\item \textbf{Dynamic dependence.} Crippa's BLM model is static. BLM's dynamic
model and LLMR's structural model are outside the stripped-down BLM equation. This
is fine for a static nesting argument, but not for claims about wage dynamics, previous
employer effects, or endogenous mobility.
\end{enumerate}

Bottom line: Crippa's stripped-down BLM is the right object if we want the clean
functional-form comparison
\[
\text{TWFE} \subset \text{Tukey} \subset \text{firm-specific-slope BLM}.
\]
It is not the full empirical BLM framework. The main thing he drops that could matter
for us is BLM's dimension reduction. Without grouping, the graph requirement becomes
so strong that identification may fail for a large share of real firms. That is not a
logical problem for theory; it is a warning for empirical scope.

\subsection{Brief Note on LLMR / Lamadon}

LLMR is dynamic and structural. It builds an equilibrium labor-market model with
worker heterogeneity, firm heterogeneity, compensating differentials, search
frictions, and Becker-style complementarities. Source: LLMR abstract, p. 1. Workers
and firms are forward-looking; employed and unemployed workers search; workers draw
mobility shocks; matches have production $f_{xy}$ and disutility/amenity $c_{xy}$;
and wages/contracts are determined in a search-and-bargaining environment. Source:
LLMR, pp. 2--9.

For our current paper, most of LLMR is probably out of scope because it is not just a
static wage equation. It is an equilibrium model with search frictions, amenities,
endogenous mobility, surplus equations, vacancy values, and contract determination.
Source: LLMR, pp. 2--9. Importing that machinery would change the project from a
static identification/nesting exercise into a structural search model.

The reusable part is the BLM first stage. LLMR explicitly say that their
identification first uses the nonlinear estimator of BLM to recover worker and firm
types, wages, mobility, and allocations in a model-consistent way; then they invert
those objects to recover structural primitives. Source: LLMR, pp. 1--4. Thus, for our
purposes, LLMR reinforces the view that BLM is useful as a statistical/type-recovery
layer. The structural inversion from BLM outputs into production, amenities, surplus,
and policy counterfactuals is the part that is likely beyond our scope.

\paragraph{Boundary check from LLMR's identification argument.}
LLMR is not merely a bigger static wage equation. Their first-stage identified objects
include the match distribution, moving probabilities, within-job wage dynamics, and
post-move wage distributions. Source: LLMR, p. 18. Lemma 2 identifies viability,
match surplus $S_{xy}$, the worker surplus function $R_{xy}(w)$, and the wage function
$w_{xy}(R)$ using the match distribution, within-job wage-increase probabilities, and
EU/JJ mobility probabilities. Source: LLMR, pp. 18--19. Lemma 3 then identifies the
job-to-job mobility shock distribution $G^1$ from the distribution of wages after
job-to-job moves, together with $R_{xy}(w)$ and $S_{xy}$. Source: LLMR, p. 20. This
is the precise sense in which LLMR forces an enlarged interface: the sufficient data
object is not just a cross-sectional wage surface, but a mobility-conditional dynamic
wage-and-transition system.

\subsection{Takeaway for Our Nesting Claim}

It is reasonable to say our factor structure is aiming at the same reusable core as
BLM: allow worker heterogeneity and firm heterogeneity to interact, rather than
forcing additivity. But we should be precise about the level:
\[
\theta_{ij}=a_j+b_j\alpha_i
\]
nests a firm-specific-slope version of BLM's interactive mean equation. It does not by
itself nest the full BLM distributional estimator or LLMR's dynamic structural model.
If our factor form can represent $a_j+b_j\alpha_i$ as a low-rank worker-firm
interaction, then we can claim it nests the static BLM wage-function core, while
leaving BLM's grouping/mixture estimation and LLMR's structural dynamics as
implementation or interpretation layers.

\section{Modeling Section: The Low-Rank Interface}

\paragraph{Source convention.}
This section uses the synthesized note \emph{Beyond AKM note} as the project object
to be checked. Page citations to that note refer to PDF page numbers in
\texttt{loo\_full\_note\_v1.pdf}. Page citations to Shimer--Smith refer to printed
Econometrica pages. Page citations to GKLP refer to the NBER working-paper PDF
pages of Gibbons--Katz--Lemieux--Parent (2002), NBER Working Paper 8889. Page
citations to LLMR refer to PDF pages of the December 30, 2025 draft.

\subsection{Notation}

\begin{center}
\begin{tabular}{>{\raggedright\arraybackslash}p{0.24\linewidth}
                >{\raggedright\arraybackslash}p{0.68\linewidth}}
\toprule
Symbol & Meaning \\
\midrule
$Y_{ij}$ & Wage or log-wage observation on worker--firm edge $(i,j)$ in the proposed statistical interface. Source: note, p. 8. \\
$X_{ij}'\delta$ & Observed covariate component. In GKLP, this must be rich enough to include sector-specific observed-skill terms such as $X_{it}'\beta_j$ if those are not being treated as part of the interaction. Source: note, p. 8; GKLP, pp. 12--14. \\
$\alpha_i,\psi_j$ & Additive worker and firm components in the interface. Source: note, p. 8. \\
$u_i,v_j$ & Worker-side and firm-side interaction attributes, kept separate from $(\alpha_i,\psi_j)$ except in the Tukey/Crippa special case. Source: note, p. 8. \\
$\Lambda$ & Low-rank interaction matrix; $\operatorname{rank}(\Lambda)=r$. Source: note, p. 8. \\
$f(x,y)$ & Shimer--Smith flow output of a matched pair. Source: Shimer--Smith, pp. 346--347. \\
$w(x)$ & Shimer--Smith reservation wage / flow value of being unmatched, $w(x)=rW(x)$. Source: Shimer--Smith, p. 349. \\
$\pi(x\mid y)$ & Shimer--Smith endogenous flow payoff to type $x$ while matched with type $y$. Source: Shimer--Smith, p. 347. \\
$Z_i,\eta_i$ & GKLP worker components: $Z_i$ is equally valued across sectors; $\eta_i$ is differentially valued across sectors. Source: GKLP, p. 7. \\
$b_j,c_j$ & GKLP sector-specific return to $\eta_i$ and sector intercept. Source: GKLP, p. 7. \\
$m_{i,t-1}$ & GKLP posterior mean of $\eta_i$ before date $t$ under learning. Source: GKLP, pp. 8, 13. \\
$\sigma_t^2$ & GKLP posterior variance term appearing in the learning wage equation; it declines with labor-market experience and is interacted with $b_j^2$. Source: GKLP, pp. 10, 13--14. \\
$u_{it},\sigma_{it}^2$ & State-augmented notation used only for embedding a learning cross-section or experience cell in the low-rank interface when the posterior-variance object is worker-time specific. This is a derived embedding of GKLP's equation (10), not GKLP's own notation. Source for equation: GKLP, pp. 13--14. \\
\bottomrule
\end{tabular}
\end{center}

\subsection{Project Scope and the No-Hallucination Rule}

The project scope in the note is clear and, in my view, correctly disciplined:
the first paper should be a static matched-data wage-interface paper, not a full
structural search paper. The minimum viable object is
\[
Y_{ij}=X_{ij}'\delta+\alpha_i+\psi_j+u_i^\top\Lambda v_j+\varepsilon_{ij},
\qquad
\operatorname{rank}(\Lambda)=r\ll \min\{I,J\}.
\tag{$\star$}
\]
Source: note, p. 8. The note's project spine says to first fix this interface and
the reported functionals, then deliver de-biased inference for the smallest
non-additive model, and only then generalize toward rank $r$. Source: note, p. 20.

The simplification that looks right is the separation between exact nesting,
approximate low-rank projection, and full structural coverage. The exact rows are
safe only when the model's wage equation literally lands in $(\star)$ after declared
normalizations. Approximate rows require an explicit projection target and an
error-to-functional bound. The note already flags this as open. Source: note,
pp. 22, 24--25. This distinction prevents the main hallucination risk: saying that
the factor interface ``nests'' a whole dynamic structural model when it only nests a
static wage-equation slice.

With that convention, the note's broad claim is mostly right but needs careful
wording:
\[
\text{AKM/KSS mean model} \subset \text{Tukey/Crippa} \subset
\text{low-rank factor interface},
\]
and BLM's static type-pair wage surface is inside the interface after taking type
indicators as factors. Source: note, p. 8. The qualifications are:
\[
\begin{array}{ll}
\text{KSS} & \text{is inference for additive AKM, not a new wage equation,}\\
\text{BLM} & \text{is a grouped distributional estimator, not only block means,}\\
\text{LLMR} & \text{is dynamic and uses mobility-conditional wage laws.}
\end{array}
\]
Source: KSS/Kline discussion in note, pp. 9--11; BLM comparison above; LLMR,
pp. 18--20.

\subsection{Shimer--Smith: What Fits the Interface}

Shimer--Smith do not give a log-wage regression of the form $(\star)$. Their primitive
matched output is the deterministic surface $f(x,y)$, and the endogenous flow payoff
$\pi(x\mid y)$ is pinned down by Bellman equations and Nash surplus splitting. Source:
Shimer--Smith, pp. 346--349. From the earlier Bellman derivation, we have
\[
rW(x\mid y)=\pi(x\mid y)-\delta S(x\mid y),
\qquad
S(x\mid y)=W(x\mid y)-W(x),
\]
and Shimer--Smith's Nash solution gives
\[
S(x\mid y)=S(y\mid x)
=\frac{f(x,y)-rW(x)-rW(y)}{2(r+\delta)}.
\]
Source: Shimer--Smith, p. 348; derivation in the first section of this write-up.
Using $w(x)=rW(x)$, the matched flow payoff can therefore be written as
\[
\begin{aligned}
\pi(x\mid y)
&=rW(x)+ (r+\delta)S(x\mid y) \\
&=w(x)+\frac{f(x,y)-w(x)-w(y)}{2}\\
&=\frac{1}{2}f(x,y)+\frac{1}{2}w(x)-\frac{1}{2}w(y).
\end{aligned}
\tag{SS-payoff}
\]
This last display is derived from Shimer--Smith's equations (3)--(4), pp. 348--349;
it is not written in the paper in this exact form.

Equation (SS-payoff) is the cleanest way to map Shimer--Smith into the project
language. It says that the wage/payoff surface is additive plus one-half of the
production surface:
\[
\pi(x\mid y)=\alpha_x+\psi_y+c_{xy},
\qquad
\alpha_x=\frac{1}{2}w(x),\quad
\psi_y=-\frac{1}{2}w(y),\quad
c_{xy}=\frac{1}{2}f(x,y).
\]
Therefore Shimer--Smith fits $(\star)$ \emph{if and only if} the relevant production
surface $f(x,y)$, or its double-centered component after absorbing additive parts,
has low rank or is well approximated by low rank:
\[
f(x,y)=f_x+f_y+2u_x^\top\Lambda v_y
\quad\Longrightarrow\quad
\pi(x\mid y)=\tilde\alpha_x+\tilde\psi_y+u_x^\top\Lambda v_y.
\]
This implication is derived from (SS-payoff). It is not an assumption in
Shimer--Smith. Their assortative-matching assumptions are supermodularity/log-
supermodularity conditions on $f$ that discipline matching sets; they do not imply
that $f$ is low rank. Source: Shimer--Smith, pp. 344--346, 349--363. This confirms
the caution in the note's coverage table: generic Shimer--Smith is better labeled
``outside unless $f$ is low-rank or an approximation theory is supplied,'' not an exact
nested row. Source: note, pp. 24--26.

\subsection{GKLP in the Interface}

GKLP start from output
\[
y_{ijt}=\exp(X_{it}'\beta_j+\psi_{ijt}),
\]
with
\[
\psi_{ijt}=Z_i+b_j(\eta_i+\varepsilon_{ijt})+c_j.
\]
Source: GKLP, pp. 6--7, equations (1)--(2). They distinguish $Z_i$, which is equally
valued in all sectors, from $\eta_i$, which is differentially valued across sectors.
Source: GKLP, p. 7.

In the perfect-information comparative-advantage case, GKLP's log wage is
\[
\ln w_{ijt}=X_{it}'\beta_j+Z_i+b_j\eta_i+c_j+\frac{1}{2}b_j^2\sigma_\varepsilon^2.
\tag{GKLP-PI}
\]
Source: GKLP, p. 12, equation (9). If $X_{it}'\beta_j$ is included in the observed
covariate part $X_{ij}'\delta$, then (GKLP-PI) is exactly rank one:
\[
\alpha_i=Z_i,\qquad
\psi_j=c_j+\frac{1}{2}b_j^2\sigma_\varepsilon^2,\qquad
u_i=\eta_i,\qquad
v_j=b_j,\qquad
\Lambda=1.
\]
This is an exact substitution into $(\star)$, so the note's static GKLP row is correct.
Source for the row: note, p. 8; source for the wage equation: GKLP, p. 12.

Under learning and comparative advantage, GKLP's log wage is
\[
\ln w_{ijt}
=X_{it}'\beta_j+Z_i+b_jm_{i,t-1}+c_j+\frac{1}{2}b_j^2\sigma_t^2.
\tag{GKLP-L}
\]
Source: GKLP, p. 13, equation (10). For a single cross-section or experience cell,
this again has a low-rank representation once the time/experience term is handled:
$b_jm_{i,t-1}$ is rank one, and the $b_j^2\sigma_t^2/2$ term is either a firm/sector
component when $\sigma_t^2$ is common within the slice, or a second factor if one
lets the worker-side state include the relevant posterior-variance/experience object:
\[
u_{it}=
\begin{pmatrix}
m_{i,t-1}\\
\sigma_{it}^2
\end{pmatrix},
\qquad
v_j=
\begin{pmatrix}
b_j\\
\frac{1}{2}b_j^2
\end{pmatrix},
\qquad
\Lambda=I_2.
\]
The rank-two display is a derived embedding of (GKLP-L), matching the note's
Remark 1. Source: note, pp. 8--9; GKLP, pp. 13--14.

The important boundary is the pooled panel. In GKLP, $m_{i,t-1}$ is a posterior
belief updated from the worker's output history, and sector choice is endogenous to
expected ability. Source: GKLP, pp. 8, 13--16. So the pooled learning model is not a
static worker--firm equation with time-invariant $u_i$. It is covered only as repeated
cross-sectional/state-augmented slices, or by adding dynamic machinery. This is exactly
the right way to state the note's conclusion: GKLP perfect-information/static
comparative advantage is exact rank one; GKLP learning is rank one or rank two per
slice depending on how the posterior-variance term is absorbed; the full pooled panel is
at the dynamic boundary. Source: note, pp. 8--9, 21--22.

\subsection{Bottom Line for the Nesting Claim}

The nesting structure in the screenshot is basically right if the claim is about
\emph{wage-equation cores}. The safest phrasing is:
\[
\begin{array}{ll}
\text{AKM/KSS mean model:} & \Lambda=0,\\[0.25em]
\text{Tukey/Crippa:} & u_i=\alpha_i,\; v_j=\psi_j,\; \Lambda=\beta_0,\\[0.25em]
\text{BLM static type means:} & u_i=e_{k(i)},\; v_j=e_{\ell(j)},\; \Lambda=M,\\[0.25em]
\text{GKLP perfect information:} & u_i=\eta_i,\; v_j=b_j,\; \Lambda=1,\\[0.25em]
\text{Shimer--Smith restricted:} & c_{xy}=\frac{1}{2}f(x,y) \text{ is low-rank.}
\end{array}
\]
The two qualifications I would keep visible next week are:

\begin{enumerate}
\item \textbf{Do not say the interface nests full BLM, KSS, GKLP learning, or LLMR.}
It nests the static wage-equation cores. KSS contributes the leave-out inference problem;
BLM contributes grouping, mixtures, and distributional estimation; GKLP learning and
LLMR contribute dynamic state/mobility objects. Source: note, pp. 20--22, 24--26;
LLMR, pp. 18--20.

\item \textbf{Do not say generic Shimer--Smith is low rank.} The derived payoff equation is
\[
\pi(x\mid y)=\frac12 f(x,y)+\frac12 w(x)-\frac12 w(y).
\]
So the interaction is exactly $\frac12 f(x,y)$ after additive terms are absorbed.
This lands in $(\star)$ only if $f$ is low rank or explicitly approximated by low rank.
Source: Shimer--Smith, pp. 346--349; note, pp. 24--26.
\end{enumerate}

Thus, the current factor form plausibly covers the vast majority of \emph{static wage
surfaces people actually estimate or summarize}: AKM/KSS means, Crippa/Tukey,
static BLM type-cell means, and GKLP's static comparative-advantage equation. The
honest frontier is dynamic structural content: amenities, on-the-job search, mobility
histories, contracts, and policy counterfactuals. That frontier is important, not marginal,
but it should be presented as the next interface enlargement rather than as something
already nested by $(\star)$.

\section{Kline Review: Edge Effects, Inference, and the Open Margin Notes}

\paragraph{Source convention.}
Page citations to Kline refer to the PDF page numbers in the NBER working-paper
PDF, \emph{Firm Wage Effects}, October 2024, revised February 2025. The formulas
below are either Kline's formulas or algebraic translations of them. When I map the
Tukey interaction into Kline's edge-effect notation, I label the step as derived.

\subsection{Notation}

\begin{center}
\begin{tabular}{>{\raggedright\arraybackslash}p{0.24\linewidth}
                >{\raggedright\arraybackslash}p{0.68\linewidth}}
\toprule
Symbol & Meaning \\
\midrule
$Y_{it}$ & Log wage of worker $i$ in period $t$. Source: Kline, p. 15. \\
$j(i,t)$ & Firm employing worker $i$ in period $t$. Source: Kline, p. 15. \\
$\alpha_i,\psi_j$ & AKM worker and firm effects. Source: Kline, p. 15. \\
$X_{it}'\beta$ & Time-varying covariate component, including year and experience controls. Source: Kline, p. 15. \\
$R$ & Adjusted wage change $Y_2-Y_1-(X_2-X_1)'\beta$ in Kline's two-period edge notation. Source: Kline, p. 18. \\
$E$ & Worker-by-edge matrix of directed origin-destination edge indicators. Source: Kline, p. 19. \\
$B$ & Firm-by-edge incidence matrix: each edge column has $-1$ at the origin and $+1$ at the destination. Source: Kline, p. 18. \\
$\Delta$ & Unrestricted vector of directed edge effects. Source: Kline, p. 19. \\
$c$ & A vector in the cycle space of the mobility graph; equivalently, $Bc=0$. Source: Kline, p. 18. \\
$C\dot\eta$ & Cycle component in Kline's decomposition $\Delta=B'\dot\psi+C\dot\eta$. Source: Kline, p. 20. \\
$W$ & Diagonal edge-weight matrix recording the number of workers moving along each edge. Source: Kline, pp. 20, 27. \\
$\hat\Delta,\tilde\Delta$ & Estimated unrestricted edge effects and AKM-predicted edge effects. Source: Kline, pp. 24--27. \\
$H$ & Weighted projection matrix mapping estimated edge effects to AKM-predicted edge effects, $\tilde\Delta=H\hat\Delta$. Source: Kline, p. 24. \\
$M$ & Residual-maker matrix $M=I-H$. Source: Kline, p. 27. \\
$h_{\ell\ell}$ & Leverage of edge $\ell$ in the AKM edge projection. Source: Kline, pp. 24, 27. \\
$n_\ell$ & Number of movers traversing edge $\ell$. Source: Kline, p. 27. \\
$E_u[\cdot],V_u[\cdot]$ & Expectation and variance with respect to the edge-effect wage-change error $u$. Source: Kline, p. 27. \\
\bottomrule
\end{tabular}
\end{center}

\subsection{What Section 3 Establishes}

Kline starts from the AKM model
\[
Y_{it}=\alpha_i+\psi_{j(i,t)}+X_{it}'\beta+\varepsilon_{it}.
\]
Source: Kline, p. 15. His useful reminder is that the strict-exogeneity condition does
two jobs at once: it rules out mobility driven by time-varying wage shocks, and it
imposes additive separability of expected log wages in worker and firm heterogeneity.
Source: Kline, p. 15.

In the two-period mover graph, AKM implies
\[
R=EB'\psi+\varepsilon,
\]
where $R$ is the adjusted wage change and $E$ records which directed firm-to-firm
edge each mover traverses. Source: Kline, pp. 18--19. The unrestricted edge-effect
model is
\[
R=E\Delta+u.
\]
Source: Kline, p. 19. The AKM restriction is therefore
\[
\Delta=B'\psi.
\]
Source: Kline, p. 19.

The key graph-theoretic implication is that AKM sets every cycle sum to zero:
\[
c'\Delta=c'B'\psi=0
\qquad \text{for every cycle vector } c.
\]
Source: Kline, pp. 18--19. Kline also states the converse: within a connected set,
these cyclic restrictions exhaust the restrictions that AKM places on edge effects.
He writes
\[
\Delta=B'\dot\psi+C\dot\eta,
\]
where $C$ collects fundamental cycles, and AKM amounts to $\dot\eta=0$. Source:
Kline, p. 20. This is the clean connection to Crippa: cycles are where additive
firm effects become testable.

\subsection{Section 3.2.2: Accounting for Noise}

The inference issue in Section 3.2.2 is that estimated edge effects are noisy. A large
residual between $\hat\Delta$ and the AKM prediction $\tilde\Delta$ does not by itself
prove model failure, because some residual variation would arise even if AKM were
true. Kline writes the AKM prediction as
\[
\tilde\Delta=H\hat\Delta,
\qquad
\hat\Delta-\tilde\Delta=M\hat\Delta,
\qquad
M=I-H.
\]
Source: Kline, pp. 24, 27. For the mover-weighted residual sum of squares,
\[
\hat\Delta'M'WM\hat\Delta,
\]
Kline decomposes the expectation as
\[
E_u[\hat\Delta'M'WM\hat\Delta]
=\Delta'M'WM\Delta+
\operatorname{trace}\!\left(M'WMV_u[\hat\Delta]\right).
\]
Source: Kline, p. 27. The first term is model error or squared approximation error.
The second term is sampling noise in the estimated edge effects. Under exact AKM,
$M\Delta=0$, so the first term vanishes, but the trace term remains. Source: Kline,
p. 27.

Under independent wage-change errors across movers, the trace term simplifies to
\[
\sum_{\ell\in E} n_\ell(1-h_{\ell\ell})V_u[\hat\Delta_\ell].
\]
Source: Kline, p. 27. This is the inferential object I had not fully internalized before:
high-leverage edges mechanically leave less residual variation, while noisy edges leave
more residual variation. Kline estimates $V_u[\hat\Delta_\ell]$ using within-edge
wage-change dispersion for edges with at least two movers. Source: Kline, p. 27.

Empirically, among non-bridge edges with at least two movers, the residual sum of
squares is $360.51$, while the expected noise contribution under AKM is $207.58$.
The difference, $152.92$, is Kline's estimate of squared approximation error. Relative
to a noise-adjusted edge-effect sum of squares of $978.67$, this implies AKM captures
about $84\%$ of the variation in true edge effects, with an implied edge-effect / AKM
prediction correlation of $0.92$. Source: Kline, p. 28. Across edge samples, Kline's
Table 1 gives noise-adjusted $R^2$ values ranging from roughly $72\%$ to $88\%$,
and about $81\%$ when bridges are included with imputed singleton noise. Source:
Kline, pp. 28--29.

The interpretation is not that AKM is literally true. It is that, after accounting for
edge noise, AKM is a strong low-dimensional summary of directed mover wage changes,
with remaining model error concentrated in cycles. Source: Kline, p. 29.

\subsection{Unrestricted Edge Effects and Tukey/Crippa}

Short answer: unrestricted edge effects can \emph{absorb} Tukey/Crippa interactions,
but they do not \emph{identify} the Tukey structure.

Suppose the underlying wage equation is the Tukey/Crippa mean surface
\[
m_{ij}=\alpha_i+\psi_j+\beta_0\alpha_i\psi_j.
\]
For movers from firm $j$ to firm $k$, the expected wage change contributed by this
surface is
\[
\begin{aligned}
\Delta_{jk}
&=E[m_{ik}-m_{ij}\mid i:j\to k] \\
&=(\psi_k-\psi_j)\left(1+\beta_0 E[\alpha_i\mid i:j\to k]\right).
\end{aligned}
\tag{derived}
\]
Thus $\Delta_{jk}$ is flexible enough to summarize the Tukey model's implication
for each observed mover edge. But it does \emph{not} recover $\beta_0$, $\alpha_i$,
or $\psi_j$; it only records average directed wage changes. Tukey collapses to AKM
only under special composition restrictions. For example, if
$E[\alpha_i\mid i:j\to k]$ is constant across directed edges, then
\[
\Delta_{jk}=(1+\beta_0\bar\alpha)(\psi_k-\psi_j),
\]
which is just an additive firm-effect model with rescaled firm effects. In general, edge-
specific mover composition makes the cycle sums nonzero:
\[
\sum_{(j,k)\in \text{cycle}}\Delta_{jk}\neq 0.
\]
This is derived from the Tukey formula and Kline's cycle restriction $c'\Delta=0$.
Source for the edge/cycle restriction: Kline, pp. 18--20.

\subsection{Margin Note 1: $\Delta_{jk}$ as a Statistical Interface AKM Does Not Use}

The margin note should use $\Delta_{jk}$, not $\delta_{jk}$, if we follow Kline's notation.
Kline's $\Delta_{jk}$ is the unrestricted directed edge effect: the average adjusted wage
change for movers from $j$ to $k$. Source: Kline, p. 19. AKM does not estimate this
object freely. Instead, it imposes $\Delta=B'\psi$, reducing potentially $|E|$ directed
edge effects to $J-1$ independent firm effects. Source: Kline, p. 19.

This is exactly why $\Delta_{jk}$ is a useful interface for us. It can summarize any
model's implications for observed mover wage changes, including additivity, Tukey
interactions, BLM-style type interactions, match effects, or selection. But AKM only
uses the lower-dimensional projection of these edge effects onto firm differences.
When cycles contain residual variation, that residual is the part of $\Delta$ AKM leaves
on the table. Source: Kline, pp. 20, 24--29.

\subsection{Margin Note 2: Match Effects and Cycle Restrictions}

Unrestricted match effects generally violate AKM cycle restrictions. To see this, let
the wage surface include a worker-firm match component $\mu_{ij}$:
\[
m_{ij}=\alpha_i+\psi_j+\mu_{ij}.
\]
Then the directed edge effect from $j$ to $k$ is
\[
\Delta_{jk}
=\psi_k-\psi_j+
E[\mu_{ik}-\mu_{ij}\mid i:j\to k].
\tag{derived}
\]
The first term telescopes around a cycle; the second generally does not, because
different directed edges are traversed by different selected mover populations. Hence
match effects create a cycle component $C\dot\eta$ in Kline's decomposition unless
their edge-specific averages happen to satisfy the same zero-cycle restrictions as firm
effects. Source for the cycle decomposition: Kline, p. 20.

This is the cleanest way to connect the margin note to the project. Match effects are
not merely ``extra residual noise.'' If their conditional means vary systematically across
directed mover edges, they show up as violations of AKM's cycle restrictions and as
model error in Kline's noise-adjusted goodness-of-fit exercise. Source: Kline, pp. 27--29.

\subsection{Connection to Event-Study Diagnostics}

Kline's causal section interprets the AKM model as a many-treatment difference-in-
differences problem. His Assumption 2 is a parallel-trends condition: among workers
switching between any pair of firms, average potential wages at the origin firm would
not have changed between periods. Source: Kline, p. 30. He notes that the standard
event-study diagnostic from Card--Heining--Kline plots mover wage trajectories by
origin/destination firm wage quartiles and asks whether movers to high-wage firms
already had faster wage growth before the move. Source: Kline, p. 30.

This matches the Sorkin--Warwar discussion above. The event study is a diagnostic
for a parallel-trends/no-match-effect style restriction, but SW show it can have low
power against endogenous match-driven mobility because pooled event-study slopes can
look close to one even when subsamples reveal failures. Source: SW, pp. 31--35; Kline,
p. 30.

\section{Dropbox Note Cleanup: Section 10 Open Items}

\paragraph{Source convention.}
Page citations to ``the note'' refer to PDF page numbers in
\texttt{loo\_full\_note\_v1.pdf}. I treat the note as an internal project document:
when it states a result, I record whether this write-up has verified it against an
external source, derived it here, or left it as future work.

\subsection{Notation}

\begin{center}
\begin{tabular}{>{\raggedright\arraybackslash}p{0.24\linewidth}
                >{\raggedright\arraybackslash}p{0.68\linewidth}}
\toprule
Symbol & Meaning \\
\midrule
Gate 1 & Fix the interface, estimands, exact nestings, and coverage audit. Source: note, p. 20. \\
Gate 2 & First new inference theorem for the smallest non-additive model, especially de-biased functionals beyond additivity. Source: note, p. 20. \\
Gate 3 & Extension from rank one to rank $r$ after normalizations, graph conditions, and computation are settled. Source: note, p. 20. \\
T1--T3 & The note's tiers of open items: Tier 1 blocks the minimum viable paper, Tier 2 strengthens the core, Tier 3 is downstream/parked. Source: note, pp. 20--22. \\
Block W & Wage-equation/interface block. Source: note, pp. 26--33. \\
Block G & Match-graph / acceptance-probability block used for policy objects such as ITT. Source: note, pp. 31--33. \\
ITT, ATT & Intent-to-treat and average-treatment-on-treated meeting values in Borovi\v{c}kova--Shimer. Source: BS, pp. 10--11. \\
$R_{xy},S_{xy}$ & Worker surplus and maximum feasible worker surplus in LLMR; part of the enlarged dynamic interface. Source: LLMR, pp. 18--20. \\
\bottomrule
\end{tabular}
\end{center}

\subsection{Status Audit}

The note's sequencing is correct: the paper should be built around the intersection of
important reported functionals, models already used, and feasible new inference. Source:
note, p. 20. The cleanup status is:

\begin{longtable}{>{\raggedright\arraybackslash}p{0.16\linewidth}
                  >{\raggedright\arraybackslash}p{0.30\linewidth}
                  >{\raggedright\arraybackslash}p{0.42\linewidth}}
\toprule
Item & Status & Reason / citation \\
\midrule
\endfirsthead
\toprule
Item & Status & Reason / citation \\
\midrule
\endhead
T1.1 out-of-scope triage & Partly closed & The dynamic branch should be called important but outside, not marginal. LLMR verifies this: its identification uses mobility probabilities, within-job wage dynamics, and mobility-conditional wage distributions, not just wages. Source: note, pp. 21, 24--25; LLMR, pp. 18--20. \\
Unrestricted match effects & Closed for current purposes & Kline's unrestricted edge effects absorb match effects statistically, but AKM does not; match effects generally create cycle components. Source: Kline, pp. 19--20, 27--29. \\
T2.1 rank-$r$ identification & Still open theory & Crippa/Tukey and Appendix H are reviewed above, but the note's requested adversarial pass on the rectangle-graph condition, necessity, and data verification remains real theory/data work. Source: note, p. 21. \\
T2.2 BS off-Pareto rank & Partly checked, not closed & The note's claimed generalized-Pareto rank result remains an internal derived claim. BS Appendix E confirms that outside Pareto, the average log wage can depend on the whole production function $f$, and in the exponential case multiplicative $f$ gives a structured wage surface, not automatic exact additivity. Source: note, p. 21; BS, pp. 69--73. \\
T2.3 BS pseudo-true error bound & Open & This requires a formal projection target and a bound from approximation error to functional error. The note itself labels this as pending. Source: note, pp. 21--22. \\
T2.4 BS bookkeeping & Partly checked & BS provides the selection threshold and ATT/ITT formulas; the note's stronger bookkeeping claim about $(\gamma,\theta,z_0,\kappa)$ still needs a line-by-line pass before being quoted as closed. Source: note, p. 21; BS, pp. 6--11. \\
T2.5 paired-mover validation & Partly closed & Kline gives the edge-effect diagnostic and noise-adjusted model-error framework; SW gives the low-power event-study warning. The formal population object for a paired-mover regression under low-rank projection remains open. Source: note, p. 21; Kline, pp. 27--30; SW, pp. 31--35. \\
T2.6 GKLP learning & Closed for current scope & Perfect-information GKLP is exact rank one; learning is rank one/rank two by slice and dynamic at the pooled-panel level. Source: note, pp. 21--22; GKLP, pp. 12--16. \\
T2.7 novelty/literature sweep & Open & This still needs a web/literature pass on de-biased functionals combining two-way outcome and graph-formation blocks. Source: note, p. 22. \\
T3.1 BS Appendix E & Partly closed & Appendix E was checked enough to know the exponential case is not a free exact-nesting result; full reconciliation with the note's Proposition 7 should be kept as a small follow-up. Source: note, p. 22; BS, pp. 69--73. \\
T3.2 Block-G cycle scarcity & Open data task & Requires counting informative cycles in the target empirical graph. Source: note, p. 22. \\
T3.3 LLMR boundary & Partly closed & Claim (b) is supported: LLMR needs surplus, wage dynamics, and mobility-conditional wage laws. The larger class-sufficiency theorem remains parked. Source: note, pp. 22, 36; LLMR, pp. 18--20. \\
T3.4 Target B theorem & Parked & The note itself says to park it until the core gates clear. Source: note, p. 22. \\
T3.5 pseudo-true projection & Open and important & Until this is formalized, coverage should be claimed only for exact rows; approximate rows should stay labeled approximate. Source: note, p. 22. \\
\bottomrule
\end{longtable}

\subsection{Clean Boundary Map After the Cleanup}

The current paper-facing boundary should read as follows:

\begin{enumerate}
\item \textbf{Covered exactly as static wage-equation cores:} AKM/KSS mean model,
Tukey/Crippa, static BLM type-cell means, and GKLP perfect-information comparative
advantage. Source: note, p. 8; Crippa section above; GKLP, p. 12.

\item \textbf{Covered only by static slices or with explicit qualifications:} GKLP learning
and low-rank restrictions/approximations to Shimer--Smith or BS off-Pareto wage
surfaces. Source: note, pp. 8--9, 21--26; Shimer--Smith, pp. 346--349; BS, pp. 69--73.

\item \textbf{Diagnostic/statistical-interface objects, not structural models:} Kline's
unrestricted edge effects $\Delta_{jk}$ and paired-mover/event-study diagnostics. These
are useful because they reveal cycle violations and low-power failures, but they do not
by themselves identify $\beta_0$, match effects, production, amenities, or surplus.
Source: Kline, pp. 19--20, 27--30; SW, pp. 31--35.

\item \textbf{Important but outside the first paper's interface:} LLMR, on-the-job search
with amenities, unrestricted dynamic wage contracts, and joint Block-W--Block-G policy
counterfactuals. These require mobility-conditional dynamic objects beyond the
cross-sectional wage surface. Source: note, pp. 24--25, 36; LLMR, pp. 18--20.
\end{enumerate}

So the main correction to the Dropbox note is not a rejection of the factor structure. It
is a tightening of the verb ``nests.'' The factor model nests many static \emph{wage
surfaces}. It does not nest the full empirical/inference/structural packages around those
surfaces unless the extra assumptions and data objects are explicitly added.

\subsection{Audit Against the Original Task List}

\begin{longtable}{>{\raggedright\arraybackslash}p{0.24\linewidth}
                  >{\raggedright\arraybackslash}p{0.64\linewidth}}
\toprule
Original item & Where it is handled \\
\midrule
\endfirsthead
\toprule
Original item & Where it is handled \\
\midrule
\endhead
Shimer--Smith Bellman derivation and $E[dV]$ / loss-from-continuing term & Section 1 derives the unmatched and matched Bellman equations from the lecture HJB and the Poisson jump calculation, with notation table and page citations to Shimer--Smith and the slides. \\
Crippa Appendix H cycle questions & Section 2 answers the exactly-4-cycle versus longer-cycle question, isolates the product telescoping step for $\beta_0$, and separates connectedness for $\alpha,\psi$ from cycle traversal for $\beta_0$. \\
CHK one-page summary and Sorkin--Warwar event study & Sections 3--4 summarize CHK and SW, including the low-power event-study mechanism and its connection to Kline's parallel-trends framing. \\
Crippa BLM versus real BLM and LLMR & Section 5 compares Crippa's stripped-down BLM equation to BLM's grouped distributional framework and explains why LLMR is dynamic/out of scope while reusing BLM as a first-stage layer. \\
Modeling section and $\alpha+\psi+u'\Lambda v$ form & Section 6 checks the Dropbox note's factor structure, derives the Shimer--Smith payoff mapping, verifies GKLP from the NBER paper, and states the exact/static versus dynamic-boundary claims. \\
Kline Section 3, inference subsection, and margin notes & Section 7 covers Kline Section 3, especially 3.2.2, the unrestricted edge-effect object, Tukey/Crippa embedding, $\Delta_{jk}$ as a statistical interface, and match effects as cycle violations. \\
Dropbox note Section 10 cleanup & Section 8 audits the open items from the note, marking which are closed for current purposes, partly checked, still open, or parked. \\
\bottomrule
\end{longtable}

\end{document}
